\documentclass[12pt,oneside]{book}
\usepackage[utf8]{inputenc}
\usepackage[T1]{fontenc}
\usepackage[spanish,es-nodecimaldot]{babel}
\usepackage[top=2.5cm,bottom=2.5cm,left=3cm,right=2.5cm,headheight=15pt]{geometry}
\usepackage{amsmath,amssymb,mathtools}
\usepackage{graphicx,float,caption,subcaption}
\usepackage{booktabs,array,longtable,tabularx}
\usepackage{enumitem}
\usepackage{xcolor}
\usepackage[most]{tcolorbox}
\usepackage{fancyhdr}
\usepackage{ragged2e}
\usepackage[colorlinks=true,linkcolor=blue!50!black,citecolor=green!40!black,urlcolor=blue!60!black,bookmarks=true,pdftitle={Apuntes de Derecho Procesal Civil},pdfauthor={Isaac Edgar Camacho Ocampo},pdfsubject={Apuntes universitarios}]{hyperref}
\graphicspath{{./img/}{./graficos/}{./figuras/}}
\setcounter{tocdepth}{2}
\setlength{\parindent}{0pt}\setlength{\parskip}{0.5em}
\setlist[itemize]{topsep=0.3em,itemsep=0.2em}
\setlist[enumerate]{topsep=0.3em,itemsep=0.2em}
\clubpenalty=10000\widowpenalty=10000
\newcommand{\materia}{Derecho Procesal Civil}
\newcommand{\titulo}{Clase introductoria: Unidades 1 y 2}
\newcommand{\autor}{Isaac Edgar Camacho Ocampo}
\newcommand{\anio}{2026}
\pagestyle{fancy}\fancyhf{}
\fancyhead[L]{\nouppercase{\leftmark}}\fancyhead[R]{\materia}\fancyfoot[C]{\thepage}
\renewcommand{\headrulewidth}{0.4pt}\renewcommand{\footrulewidth}{0pt}
\fancypagestyle{plain}{\fancyhf{}\fancyfoot[C]{\thepage}\renewcommand{\headrulewidth}{0pt}}
\newcommand{\abs}[1]{\left|#1\right|}
\newcommand{\norm}[1]{\left\lVert#1\right\rVert}
\newcommand{\vect}[1]{\mathbf{#1}}
\newcommand{\deriv}[2]{\frac{d#1}{d#2}}
\newcommand{\pderiv}[2]{\frac{\partial#1}{\partial#2}}
\newtcolorbox{definition}[1]{enhanced,breakable,title={Definición: #1},fonttitle=\bfseries,colback=blue!3,colframe=blue!50!black,boxrule=0.8pt,arc=2mm}
\newtcolorbox{important}{enhanced,breakable,title={Importante},fonttitle=\bfseries,colback=yellow!8,colframe=orange!70!black,boxrule=0.8pt,arc=2mm}
\newtcolorbox{example}[1]{enhanced,breakable,title={Ejemplo: #1},fonttitle=\bfseries,colback=green!4,colframe=green!40!black,boxrule=0.8pt,arc=2mm}
\newtcolorbox{formula}[1]{enhanced,breakable,title={Fórmula / Regla: #1},fonttitle=\bfseries,colback=purple!4,colframe=purple!50!black,boxrule=0.8pt,arc=2mm}
\newtcolorbox{exercise}[1]{enhanced,breakable,title={Ejercicio: #1},fonttitle=\bfseries,colback=gray!5,colframe=gray!60!black,boxrule=0.8pt,arc=2mm}
\newtcolorbox{note}{enhanced,breakable,title={Nota},fonttitle=\bfseries,colback=white,colframe=gray!50,boxrule=0.6pt,arc=2mm}
\begin{document}
\begin{titlepage}\thispagestyle{empty}\begin{center}
\vspace*{1cm}
{\large\textsc{Universidad de Buenos Aires}}
\vspace{2.5cm}
{\Huge\textbf{\materia}}
\vspace{1cm}
{\LARGE\titulo}
\vspace{0.8cm}
\textit{Comprender el proceso permite transformar el estudio jurídico en una herramienta para la práctica.}
\vspace{1.2cm}\rule{0.70\textwidth}{0.5pt}\vspace{0.5cm}
{\large Apuntes de estudio}
\vspace{0.5cm}\rule{0.70\textwidth}{0.5pt}
\vfill{\large\autor}\vspace{0.3cm}{\large\anio}
\end{center}\vfill
\begin{flushleft}\textit{Este es un modesto aporte a los estudiantes de ``Derecho Procesal Civil''; no pretende sustituir el material oficial de la materia.}\end{flushleft}
\end{titlepage}
\tableofcontents
\chapter{Presentación y modalidad de cursada}
Las clases son grabadas y se suben al campus virtual de la facultad y a YouTube, lo que permite consultarlas en cualquier momento. La estructura de seguimiento incluye profesores adjuntos, jefes de trabajos prácticos y ayudantes. Cada comisión cuenta, además, con una página de Facebook como canal de comunicación.

\section{Clases prácticas y evaluación}
Las clases prácticas permiten observar cómo se realiza una demanda, cómo se contesta y cómo se desarrolla un interrogatorio de testigos. Esto resulta especialmente relevante porque la materia puede presentarse como abstracta para los estudiantes ingresantes.

El sistema de evaluación contempla un primer parcial, un segundo parcial y sus respectivos recuperatorios. El examen es con libro abierto y permite utilizar el Código Procesal. Esta modalidad no elimina la dificultad: exige comprensión y no una mera memorización.

\begin{important}
El estudio del Derecho Procesal debe orientarse a la comprensión. La materia constituye una herramienta para el trabajo práctico cotidiano del abogado.
\end{important}

\section{Material de estudio}
Se menciona el libro de la cátedra, \textit{Nociones Básicas del Derecho Procesal Civil}, en prensa. Mientras tanto, se recomiendan obras de Palacio, Falcón y Arazi, junto con el Código Procesal Civil y Comercial y sus leyes complementarias.

\chapter{El Derecho Procesal Civil y el conflicto jurídico}
\section{Derecho sustancial y derecho procesal}
El derecho sustancial o derecho de fondo regula la vida ordinaria en sociedad. El derecho procesal interviene cuando aparece un conflicto y las normas sustanciales dejan de ser operativas entre las partes.

\begin{definition}{Derecho Procesal Civil}
Es la herramienta que permite hacer actuar la voluntad de la ley cuando el derecho de fondo deja de ser operativo entre las partes debido a la aparición de un conflicto.
\end{definition}

Cuando el conflicto se judicializa aparece un tercero imparcial: el juez. La relación jurídica que se desarrolla a partir de ese momento se rige por el Código Procesal Civil y Comercial. La sentencia constituye una norma individual, una ``ley a medida'' del caso concreto, que vuelve operativa la voluntad del derecho sustancial por mandato judicial.

\section{Ejemplos del surgimiento del proceso}
\begin{example}{Accidente de colectivo}
Una persona celebra un contrato de transporte, el colectivo choca y el pasajero sufre una lesión que requiere una semana de internación. Si la empresa no responde, la relación sustancial se vuelve conflictiva y surge una nueva relación jurídica en la que interviene el juez.
\end{example}
\begin{example}{Compra de una heladera}
Una persona compra una heladera que no funciona, la garantía no responde y el local cierra. La relación de consumo queda insatisfecha y puede acudirse al Tribunal de Consumo o a la justicia ordinaria mediante el proceso.
\end{example}
\begin{example}{Infortunio laboral}
Un trabajador manipula una máquina, una pieza rota provoca una caída y sufre una conmoción cerebral. La relación laboral sustancial da lugar a un conflicto que debe canalizarse por la vía procesal correspondiente.
\end{example}

\begin{formula}{Normas aplicables}
El Código Civil y Comercial de la Nación rige la relación sustancial; el Código Procesal Civil y Comercial rige el desarrollo del proceso una vez surgido el conflicto.
\end{formula}

\chapter{Autonomía, unidad y fuentes}
\section{Autonomía y unidad del Derecho Procesal}
El Derecho Procesal es una materia autónoma e independiente. Posee un código propio, doctrina propia y autores especializados. Aunque existan ramas especializadas, la cátedra sostiene que el derecho procesal es uno solo, porque todas sus manifestaciones se apoyan en un mismo pedestal básico: \textbf{jurisdicción, acción o pretensión y proceso}.

\subsection{Ramas especializadas}
\begin{itemize}
\item \textbf{Derecho Procesal Penal:} incluye, entre otras figuras, el Ministerio Público como parte acusadora, recursos de nulidad, casación e inconstitucionalidad y medidas cautelares como la prisión preventiva.
\item \textbf{Derecho Procesal Laboral:} incorpora el principio protectorio, que busca equilibrar la asimetría entre trabajador y empleador.
\item \textbf{Derecho Procesal Constitucional:} estudia el amparo, la acción declarativa de inconstitucionalidad y el recurso extraordinario.
\item \textbf{Derecho Procesal Administrativo:} se vincula con el control jurisdiccional de los actos administrativos.
\end{itemize}

\subsection{Uso práctico del Código}
Se recomienda tener siempre a mano un Código Procesal comentado, comparado con el estetoscopio del médico. La edición denominada ``gordita'' incluye leyes complementarias, entre ellas la Ley 22.172 sobre exhortos entre jurisdicciones del país.

\section{Fuentes del Derecho Procesal}
\begin{enumerate}
\item Constitución Nacional, especialmente el artículo 18 y las disposiciones orgánicas relativas al Poder Judicial.
\item Tratados internacionales de derechos humanos, vinculados con el artículo 75, incisos 22, 23 y 24 de la Constitución Nacional.
\item Leyes nacionales y locales.
\item Decretos, incluidos los decretos de necesidad y urgencia. La Ley 26.122 regula su control parlamentario mediante la Comisión Bicameral.
\item Jurisprudencia, especialmente los fallos de la Corte Suprema de Justicia de la Nación.
\item Doctrina elaborada por autores especializados.
\end{enumerate}

Se mencionan el caso \textit{Ekmekdjian c/ Sofovich}, relativo a la operatividad de derechos consagrados en tratados internacionales, y los DNU dictados durante la pandemia, que incidieron en el derecho procesal al suspender lanzamientos y ejecuciones en el proceso laboral.

\section{Dinámica de la norma procesal}
Una misma situación puede producir consecuencias procesales diferentes. Por ejemplo, una vez notificada la demanda y trabada la litis, el demandado puede comparecer y contestar, comparecer sin contestar o no comparecer. Cada alternativa genera consecuencias reguladas por el Código.

\chapter{Breve historia del Derecho Procesal}
\section{Antecedentes históricos}
El proceso tiene raíces en el procedimiento oral romano, desarrollado en dos etapas: \textit{in iure}, ante un magistrado, y \textit{apud iudicem}, ante un juez o árbitro elegido por las partes.

Con las invasiones bárbaras, el derecho romano recibió la influencia del derecho germánico, que incorporaba elementos de carácter divino, como las pruebas de ordalías. Durante la Edad Media, la influencia de la Iglesia dio lugar al llamado derecho común, que llegó a la Argentina principalmente a través de España y de las Leyes de Enjuiciamiento Civil de 1855 y 1881.

\section{Sistemas de justicia mencionados}
\begin{itemize}
\item \textbf{Sistema de la tribu Azande:} se utilizaba un procedimiento basado en la reacción de un pollo ante la inoculación de doble veneno.
\item \textbf{Sistema anglosajón (\textit{common law}):} el jurado y la declaración testimonial ocupan un lugar central; la prueba testimonial posee mayor peso que la documental.
\item \textbf{Sistema continental (\textit{civil law}):} sistema al que pertenece Argentina, caracterizado en el material por la preponderancia de la prueba documental y del expediente escrito.
\end{itemize}

\section{Código Procesal de 1967}
El Código Procesal Civil y Comercial argentino fue sancionado en 1967, con fuerte influencia del derecho procesal italiano. Desde entonces experimentó numerosas modificaciones. El material lo caracteriza como un código escrito y formal, con predominio histórico del expediente en papel, aunque en proceso de incorporación de oralidad y digitalización.

\chapter{El pedestal básico del Derecho Procesal}
\section{Jurisdicción}
\begin{definition}{Jurisdicción}
Potestad del Poder Judicial de resolver conflictos mediante una norma individual, es decir, de ``decir el derecho'' (\textit{iuris dictio}).
\end{definition}

\subsection{Jurisdicción y competencia}
La jurisdicción es una función constitucional del Poder Judicial. La competencia es la organización del trabajo de los tribunales, establecida por el legislador, y delimita el ámbito concreto de actuación de la jurisdicción.

La competencia puede determinarse por:
\begin{itemize}
\item territorio;
\item materia;
\item turno;
\item grado.
\end{itemize}
Se mencionan los artículos 116 y 117 de la Constitución Nacional en relación con la competencia apelada y la competencia originaria y exclusiva de la Corte Suprema.

\subsection{Momentos de la jurisdicción}
\begin{enumerate}
\item \textbf{Notio:} facultad de conocer el problema y los hechos de la demanda.
\item \textbf{Vocatio:} facultad de convocar a las partes al proceso.
\item \textbf{Coertio:} poder de coerción para asegurar el desarrollo del proceso.
\item \textbf{Iudicium:} dictado de la sentencia, entendida como norma individual.
\item \textbf{Executio:} potestad de ejecutar lo decidido cuando no existe cumplimiento voluntario.
\end{enumerate}

\begin{example}{Executio}
Si se condena a una persona a escriturar un inmueble y no lo hace, el juez puede suscribir la escritura en su nombre.
\end{example}

\subsection{Imperio y jurisdicción}
El imperio es la potestad de utilizar la fuerza pública para hacer cumplir decisiones. No es exclusivo del Poder Judicial: también lo poseen el Poder Ejecutivo y el Poder Legislativo en sus respectivos ámbitos. La jurisdicción, entendida como potestad de decir el derecho, se presenta como función privativa del Poder Judicial, con los matices indicados para ciertos procedimientos no judiciales sujetos a revisión.

\subsection{Manifestaciones no judiciales}
\begin{itemize}
\item \textbf{Jurisdicción administrativa:} la Administración puede tramitar procedimientos y aplicar sanciones, pero debe existir posibilidad de control judicial. Se menciona \textit{Fernández Arias c/ Poggio}.
\item \textbf{Jurisdicción eclesiástica:} la Sacra Rota Romana puede tramitar la nulidad de un matrimonio canónico dentro del ámbito eclesiástico.
\item \textbf{Entes reguladores:} ejercen una jurisdicción limitada o ``jurisdicción primaria''. Se menciona el caso \textit{Ángel Estrada y Cía.}
\item \textbf{Jurisdicción arbitral:} los árbitros son jueces privados elegidos voluntariamente por las partes. Poseen jurisdicción, pero no imperio; para ejecutar sus laudos deben recurrir a la jurisdicción oficial.
\end{itemize}

\section{Acción y pretensión}
\begin{definition}{Acción}
En sentido amplio, es el derecho constitucional de petición ante cualquier autoridad, relacionado en el material con el artículo 14 de la Constitución Nacional. Se trata de un concepto abstracto.
\end{definition}

\subsection{Polémica Windscheid--Muther}
La controversia de 1856 constituye, según el material, un hito fundacional del Derecho Procesal como ciencia autónoma. Permitió separar la acción del derecho sustancial: una persona puede accionar aun cuando finalmente se determine que no tenía razón en el fondo.

Bülow, en 1868, sostuvo que el proceso constituye una nueva relación jurídica, distinta de la relación que le dio origen. La tesis fue completada por Chiovenda en 1913.

\subsection{Acción concreta, pretensión y demanda}
Cuando la acción se ejerce concretamente ante la jurisdicción, se denomina \textbf{pretensión}. Es una manifestación de voluntad dirigida al juez contra otra parte, con el objeto de que se cumpla una determinada conducta.

\begin{important}
La relación básica puede representarse como: \textbf{actor} que pretende $\rightarrow$ \textbf{juez} como tercero imparcial $\leftarrow$ \textbf{demandado} que resiste.
\end{important}

La \textbf{demanda} es el acto formal de petición mediante el cual se inicia el proceso y que contiene la pretensión del actor. Entre los ejemplos mencionados se encuentran el reclamo de una indemnización por un accidente de colectivo, el divorcio, la indemnización por despido, la rendición de cuentas y el desalojo.

\section{Proceso}
\begin{definition}{Proceso}
Nueva relación jurídica que surge cuando un conflicto es llevado ante la jurisdicción. Se desarrolla sobre la base del Código Procesal y concluye con la sentencia.
\end{definition}

\subsection{El proceso como sistema}
El proceso puede analizarse como una relación entre actor, juez y demandado, pero también como un sistema complejo, siguiendo la Teoría General de Sistemas de von Bertalanffy. Un sistema es un conjunto de partes interrelacionadas que persiguen un objetivo común.

\subsection{Componentes del proceso}
\begin{description}
\item[Insumos:] actores, demandados, juez, empleados y funcionarios judiciales, tecnología, edificio, recursos económicos, papel y sistemas informáticos.
\item[Procesador:] etapas del proceso: fase introductoria, probatoria, conclusional y sentencia.
\item[Resultado:] la sentencia, entendida como norma individual que resuelve el conflicto.
\item[Retroalimentación:] mecanismo de control que permite corregir el resultado mediante recursos procesales, como la apelación.
\end{description}

\subsection{Etapas del proceso}
\begin{enumerate}
\item \textbf{Etapa introductoria:} demanda y contestación.
\item \textbf{Etapa probatoria:} ofrecimiento, admisión y producción de prueba.
\item \textbf{Etapa conclusional:} alegatos de bien probado.
\item \textbf{Sentencia.}
\end{enumerate}

\begin{example}{Retroalimentación}
Si una Cámara revoca la sentencia de un juez, se produce un control del resultado. El juez puede advertir que aplicó incorrectamente el derecho o valoró mal los hechos, y considerar ese criterio en un caso futuro similar.
\end{example}

\chapter{Los tres poderes del Estado y la función jurisdiccional}
El poder del Estado es uno solo y el soberano es el pueblo. Ese poder se organiza en tres departamentos con funciones diferenciadas:
\begin{itemize}
\item \textbf{Poder Legislativo:} hace la ley.
\item \textbf{Poder Ejecutivo:} administra.
\item \textbf{Poder Judicial:} dirime conflictos haciendo actuar la voluntad de la ley.
\end{itemize}

Las funciones pueden mezclarse en la práctica. El Ejecutivo legisla mediante DNU; el Legislativo desarrolla funciones administrativas y jurisdiccionales, como el enjuiciamiento de sus miembros; y el Judicial dicta acordadas de aplicación obligatoria, por ejemplo, sobre notificación electrónica. No obstante, la función básica del Poder Judicial es la jurisdicción en sentido estricto.

El material menciona un debate sobre una reforma judicial. Según lo expuesto, una reasignación de competencias dentro del fuero federal no equivaldría necesariamente a una reforma del sistema procesal: esta última requeriría modificar el sistema de procedimiento y no solamente la distribución de causas.

\chapter{Resumen Cornell: tabla de estudio}
\renewcommand{\arraystretch}{1.35}
\begin{longtable}{p{0.27\textwidth}p{0.67\textwidth}}
\toprule
\textbf{Preguntas / palabras clave} & \textbf{Notas / respuestas}\\
\midrule
\textbf{¿Qué es el Derecho Procesal y para qué sirve?} & Permite hacer actuar la voluntad de la ley cuando el derecho de fondo deja de ser operativo por la aparición de un conflicto. Es una herramienta para la práctica profesional del abogado.\\
\midrule
\textbf{¿Cuándo aparece el proceso?} & Cuando la relación jurídica sustancial se vuelve conflictiva y se necesita la intervención de un tercero imparcial: el juez.\\
\midrule
\textbf{¿Qué es la sentencia?} & Es la norma individual dictada por el juez para resolver el conflicto concreto y hacer operativa la voluntad de la ley sustancial.\\
\midrule
\textbf{¿El Derecho Procesal es uno o múltiple?} & Es uno, aunque posee especialidades como el derecho procesal penal, laboral, constitucional y administrativo. Todas se apoyan en jurisdicción, acción o pretensión y proceso.\\
\midrule
\textbf{¿Cuáles son sus fuentes?} & Constitución Nacional, tratados internacionales, leyes, decretos y DNU, jurisprudencia y doctrina.\\
\midrule
\textbf{¿Qué es la jurisdicción?} & Es la potestad de decir el derecho y resolver conflictos mediante una norma individual. Sus momentos son notio, vocatio, coertio, iudicium y executio.\\
\midrule
\textbf{¿Jurisdicción y competencia son lo mismo?} & No. La jurisdicción es una función constitucional; la competencia organiza el trabajo de los tribunales y delimita el ámbito de actuación por territorio, materia, turno y grado.\\
\midrule
\textbf{¿Qué es el imperio?} & Es la potestad de usar la fuerza pública. No es exclusiva del Poder Judicial.\\
\midrule
\textbf{¿Qué es la acción?} & Es el derecho abstracto de peticionar ante cualquier autoridad. Cuando se ejerce concretamente ante la jurisdicción se denomina pretensión.\\
\midrule
\textbf{¿Qué es la demanda?} & Es el acto formal de petición que inicia el proceso y contiene la pretensión del actor.\\
\midrule
\textbf{¿Qué es el proceso como sistema?} & Es un conjunto de partes interrelacionadas: posee insumos, etapas que funcionan como procesador, un resultado —la sentencia— y retroalimentación mediante recursos.\\
\midrule
\multicolumn{2}{p{0.94\textwidth}}{\textbf{Resumen:} El Derecho Procesal Civil permite canalizar los conflictos que vuelven inoperante al derecho sustancial. Su pedestal básico está compuesto por jurisdicción, acción o pretensión y proceso. La jurisdicción se despliega en cinco momentos; la pretensión se incorpora a la demanda; y el proceso, entendido como relación jurídica y como sistema, culmina en una sentencia que resuelve el caso concreto.}\\
\bottomrule
\end{longtable}
\end{document}
